\documentclass[11pt,a4paper]{article}

% ============================================
% PACKAGES
% ============================================
\usepackage[utf8]{inputenc}
\usepackage[T1]{fontenc}
\usepackage[english]{babel}
\usepackage{geometry}
\geometry{margin=2.5cm}
\usepackage{graphicx}
\usepackage{booktabs}
\usepackage{tabularx}
\usepackage{longtable}
\usepackage{multirow}
\usepackage{enumitem}
\usepackage{xcolor}
\usepackage{listings}
\usepackage{amsmath,amssymb}
\usepackage{hyperref}
\usepackage{cleveref}
\usepackage{caption}
\usepackage{subcaption}
\usepackage{float}
\usepackage{parskip}
\usepackage{microtype}
\usepackage{fancyhdr}

\hypersetup{
    colorlinks=true,
    linkcolor=blue!60!black,
    citecolor=green!50!black,
    urlcolor=blue!70!black
}

% ============================================
% LISTING STYLES
% ============================================

% Base style
\lstdefinestyle{base}{
    basicstyle=\ttfamily\small,
    breaklines=true,
    breakatwhitespace=true,
    columns=flexible,
    keepspaces=true,
    showstringspaces=false,
    tabsize=2,
    frame=single,
    framerule=0.4pt,
    rulecolor=\color{gray!50},
    backgroundcolor=\color{gray!3},
    xleftmargin=4pt,
    xrightmargin=4pt,
    aboveskip=8pt,
    belowskip=8pt,
    numberstyle=\tiny\color{gray},
    numbers=left,
    numbersep=6pt,
}

% JjTL
\lstdefinelanguage{jjtl}{
    morekeywords={transformation, from, to, when, helper, alert, notify, prompt, input},
    morekeywords=[2]{forall, exists, implies, such, that, with, do, if, then, else, and, or, not, is, true, false, null, in},
    sensitive=true,
    morecomment=[l]{\#},
    morecomment=[l]{--},
    morestring=[b]",
    morestring=[b]',
}

\lstdefinestyle{jjtl}{
    style=base,
    language=jjtl,
    keywordstyle=\bfseries\color{blue!70!black},
    keywordstyle=[2]\bfseries\color{purple!70!black},
    commentstyle=\itshape\color{green!40!black},
    stringstyle=\color{red!60!black},
}

% JjEL
\lstdefinelanguage{jjel}{
    morekeywords={forall, exists, implies, such, that, with, do, if, then, else, and, or, not, is, true, false, null, in},
    sensitive=true,
    morecomment=[l]{--},
    morestring=[b]",
}

\lstdefinestyle{jjel}{
    style=base,
    language=jjel,
    keywordstyle=\bfseries\color{purple!70!black},
    commentstyle=\itshape\color{green!40!black},
    stringstyle=\color{red!60!black},
}

% ATL
\lstdefinelanguage{atl}{
    morekeywords={module, create, rule, from, to, helper, context, def, using, lazy, unique, entrypoint, endpoint, do, let, in},
    morekeywords=[2]{if, then, else, endif, not, and, or, implies, self, thisModule, true, false, OclUndefined},
    sensitive=true,
    morecomment=[l]{--},
    morestring=[b]',
}

\lstdefinestyle{atl}{
    style=base,
    language=atl,
    keywordstyle=\bfseries\color{orange!70!black},
    keywordstyle=[2]\bfseries\color{purple!70!black},
    commentstyle=\itshape\color{green!40!black},
    stringstyle=\color{red!60!black},
}

% ETL (Epsilon)
\lstdefinelanguage{etl}{
    morekeywords={rule, transform, to, extends, guard, pre, post, operation, var, new, delete, return, throw},
    morekeywords=[2]{if, else, for, in, while, switch, case, default, continue, break, and, or, not, implies, self, true, false, null},
    sensitive=true,
    morecomment=[l]{//},
    morecomment=[s]{/*}{*/},
    morestring=[b]",
    morestring=[b]',
    moredelim=[is][\color{gray}]{@}{\ },
}

\lstdefinestyle{etl}{
    style=base,
    language=etl,
    keywordstyle=\bfseries\color{teal!80!black},
    keywordstyle=[2]\bfseries\color{purple!70!black},
    commentstyle=\itshape\color{green!40!black},
    stringstyle=\color{red!60!black},
}

% QVT-R
\lstdefinelanguage{qvtr}{
    morekeywords={transformation, top, relation, checkonly, enforce, domain, when, where, key, query, primitive},
    morekeywords=[2]{if, then, else, endif, not, and, or, implies, self, true, false, null, Set, Sequence, Bag, OrderedSet},
    sensitive=true,
    morecomment=[l]{--},
    morestring=[b]',
}

\lstdefinestyle{qvtr}{
    style=base,
    language=qvtr,
    keywordstyle=\bfseries\color{violet!80!black},
    keywordstyle=[2]\bfseries\color{purple!70!black},
    commentstyle=\itshape\color{green!40!black},
    stringstyle=\color{red!60!black},
}

% QVT-O
\lstdefinelanguage{qvto}{
    morekeywords={modeltype, uses, transformation, main, mapping, query, constructor, helper, init, end, object, inout, result, intermediate, class, property},
    morekeywords=[2]{if, then, else, endif, not, and, or, implies, self, true, false, null, in, var, return, while, forEach, when, where, new},
    sensitive=true,
    morecomment=[l]{--},
    morestring=[b]',
}

\lstdefinestyle{qvto}{
    style=base,
    language=qvto,
    keywordstyle=\bfseries\color{brown!80!black},
    keywordstyle=[2]\bfseries\color{purple!70!black},
    commentstyle=\itshape\color{green!40!black},
    stringstyle=\color{red!60!black},
}

% Grammar
\lstdefinestyle{grammar}{
    style=base,
    keywordstyle={},
    basicstyle=\ttfamily\small,
}

% ============================================
% CUSTOM COMMANDS
% ============================================
\newcommand{\jjel}{\textsf{JjEL}}
\newcommand{\jjtl}{\textsf{JjTL}}
\newcommand{\kw}[1]{\texttt{\textbf{#1}}}

% ============================================
% HEADER/FOOTER
% ============================================
\pagestyle{fancy}
\fancyhf{}
\fancyhead[L]{\small\textsf{JjTL \& JjEL: Language Design and Comparative Analysis}}
\fancyhead[R]{\small\thepage}
\renewcommand{\headrulewidth}{0.4pt}

% ============================================
% DOCUMENT
% ============================================

\begin{document}

% ============================================
% TITLE
% ============================================

\begin{center}
{\LARGE\bfseries JjTL and JjEL: A Set-Theoretic Approach to\\[4pt] Model Transformation and Expression Languages}

\vspace{14pt}

{\large Jjodel Project: Language Design Document}

\vspace{10pt}

{\today}
\end{center}

\vspace{6pt}

\begin{abstract}
This document presents the design and rationale of two domain-specific languages developed within the Jjodel ecosystem: \jjtl{} (Jjodel Transformation Language), a declarative Model-to-Model transformation language, and \jjel{} (Jjodel Expression Language), a set-theoretic expression language inspired by first-order logic. We describe the syntax, semantics, and, most importantly, the \emph{design motivations} behind each language construct, explaining \emph{why} certain decisions were made and what alternatives were considered. We then provide a detailed comparison with four established transformation languages: ATL, EOL/ETL, QVT Relations, and QVT Operational Mappings. Finally, we present a series of parallel transformation examples across ATL, ETL, and \jjtl{} to enable direct syntactic and conceptual comparison.
\end{abstract}

\tableofcontents
\newpage

% ============================================
% 1. INTRODUCTION
% ============================================
\section{Introduction}\label{sec:intro}

Model-Driven Engineering (MDE) relies on model transformations as a fundamental mechanism for bridging abstraction levels, synchronizing representations, and generating artifacts from high-level specifications. Over the past two decades, a rich ecosystem of transformation languages has emerged, from the OMG's standardized QVT family to community-driven tools like ATL and Epsilon. Despite this maturity, adoption of dedicated transformation languages in industrial practice remains limited, with many practitioners preferring general-purpose languages like Java or Kotlin.

We identify three root causes for this adoption gap:

\begin{enumerate}[label=(\roman*)]
    \item \textbf{Cognitive overhead}: Existing languages require learning OCL, understanding complex execution semantics (e.g., ATL's two-phase execution, QVT-R's bidirectional pattern matching), and dealing with Eclipse-specific tooling.
    \item \textbf{Syntax-semantics mismatch}: Languages like QVT-R use declarative syntax but require procedural reasoning about enforcement behavior; ATL mixes declarative bindings with imperative \kw{do} blocks under different execution models.
    \item \textbf{Expression language limitations}: OCL, while mathematically rigorous, was designed for \emph{constraint specification} (what must be true), not for \emph{value computation} (which elements, which values), yet it serves as the expression sub-language for ATL, QVT-R, and QVT-O.
\end{enumerate}

The Jjodel project addresses these issues with two complementary languages:

\begin{itemize}
    \item \textbf{\jjel{}}, a set-theoretic expression language whose core construct, \kw{forall}, has \emph{set comprehension} semantics (it returns a set, not a boolean). This bridges the gap between formal logic and practical computation, providing a single primitive that replaces \texttt{filter}, \texttt{map}, and logical quantification.
    \item \textbf{\jjtl{}}, a declarative M2M transformation language that uses \jjel{} as its expression sub-language, offering a simpler rule structure than ATL while maintaining automatic traceability and guard conditions.
\end{itemize}

\subsection{Design Criteria}

Both languages were designed under three guiding principles, applied in strict priority order:

\begin{enumerate}
    \item \textbf{Maximum declarativity.} The user describes \emph{what} they want, never \emph{how} to compute it. No assignment, no loops, no mutation.
    \item \textbf{Minimum cognitive load.} Few rules to remember, minimal syntax, no boilerplate. The most frequent MDE operations are available as native methods.
    \item \textbf{Maximum intuitiveness.} Constructs should read almost like natural language. Textual operators (\kw{and}, \kw{or}, \kw{not}, \kw{implies}), first-order-logic-inspired quantifiers (\kw{forall}, \kw{exists}), and null-safe navigation (\texttt{?.}, \texttt{??}).
\end{enumerate}

\subsection{Document Structure}

\Cref{sec:jjel} presents \jjel{} in detail, including its set-theoretic semantics and design motivations. \Cref{sec:jjtl} describes \jjtl{} and its transformation model. \Cref{sec:comparison} provides a detailed comparative analysis with ATL, EOL/ETL, QVT-R, and QVT-O. \Cref{sec:examples} presents parallel transformation examples across multiple languages. \Cref{sec:discussion} discusses open points and future directions.


% ============================================
% 2. JjEL
% ============================================
\section{JjEL: Jjodel Expression Language}\label{sec:jjel}

\jjel{} is a pure, side-effect-free expression language designed for navigating, querying, and computing values over model graphs. Every \jjel{} expression produces a value; there are no statements.

\subsection{Positioning and Scope}

\jjel{} occupies a specific niche in the MDE language landscape:

\begin{center}
\begin{tabular}{@{}lccc@{}}
\toprule
\textbf{OCL} & \textbf{\jjel{}} & \textbf{EOL} & \textbf{JavaScript} \\
\midrule
Formal constraints & Set-theoretic expressions & Model scripting & General-purpose \\
``what must be true'' & ``which elements, which values'' & ``manipulate the model'' & ``computation'' \\
\bottomrule
\end{tabular}
\end{center}

\jjel{} is \emph{not} a constraint language (it computes values, not assertions), not a programming language (no variables, loops, or I/O), and not a transformation language (it provides expressions \emph{inside} rules, not rules themselves). \jjel{} provides the expression sub-language for \jjtl{}, analogous to how OCL provides expressions for ATL and QVT.

The scope boundary is precise: \jjtl{} handles Model-to-Model (M2M) transformations; Model-to-Text (M2T) generation is handled by a separate template engine.

\subsection{Types and Values}

\jjel{} supports five primitive types and three composite types:

\begin{center}
\begin{tabular}{@{}llll@{}}
\toprule
\textbf{Category} & \textbf{Type} & \textbf{Literals} & \textbf{Examples} \\
\midrule
\multirow{5}{*}{Primitive} & \texttt{EString} & Double-quoted & \texttt{"hello"}, \texttt{"Customer"} \\
& \texttt{EInt} & Integer & \texttt{0}, \texttt{42}, \texttt{-7} \\
& \texttt{EDouble} & Decimal & \texttt{3.14}, \texttt{.5} \\
& \texttt{EBoolean} & Keyword & \texttt{true}, \texttt{false} \\
& \texttt{null} & Keyword & \texttt{null} \\
\midrule
\multirow{3}{*}{Composite} & Array & Bracket & \texttt{[]}, \texttt{[1, 2, 3]} \\
& Object & (no literal) & Model elements \\
& Function & Lambda & \texttt{a => a.name} \\
\bottomrule
\end{tabular}
\end{center}

The type system includes aliases for interoperability: \texttt{EString}/\texttt{String}, \texttt{EInt}/\texttt{Integer}, \texttt{EDouble}/\texttt{Double}/\texttt{Number}, \texttt{EBoolean}/\texttt{Boolean}. The \kw{is} operator checks type membership including supertypes (kind-of semantics).


\subsection{The Core Construct: \texttt{forall}}\label{sec:forall}

The most important design decision in \jjel{} is the semantics of \kw{forall}.

\subsubsection{Design Motivation}

In classical first-order logic, the universal quantifier $\forall$ is an \emph{assertion}: $\forall x \in S : P(x)$ returns a boolean (true if $P$ holds for every $x$). In set theory, the set-builder notation $\{x \in S \mid P(x)\}$ \emph{selects} elements, returning a set.

We chose \textbf{set-theoretic semantics}: \kw{forall} returns a \emph{set}, not a boolean. The rationale:

\begin{enumerate}
    \item \textbf{Computational utility}: In model transformations, we almost always need to \emph{select} elements and \emph{compute} values from them. A set-returning construct is directly useful; a boolean-returning one requires additional machinery.
    \item \textbf{Unification}: With set-theoretic \kw{forall}, the \texttt{filter} and \texttt{map} collection methods become redundant; \kw{forall} subsumes both. This reduces the number of concepts a user must learn.
    \item \textbf{Natural reading}: ``for all attributes such that they are public: get their name'' reads naturally and maps directly to the syntax.
\end{enumerate}

\subsubsection{Syntax and Forms}

\kw{forall} has four forms, distinguished by three separators with distinct roles:

\begin{lstlisting}[style=jjel,numbers=none]
forall <var> in <collection> [such that <condition>] [: <projection>] [do <action>]
\end{lstlisting}

\begin{center}
\begin{tabular}{@{}llll@{}}
\toprule
\textbf{Form} & \textbf{Returns} & \textbf{Equivalent to} & \textbf{Example} \\
\midrule
\texttt{forall x in S such that P} & Set & \texttt{S.filter(P)} & \texttt{forall a in attrs such that a.isPublic} \\
\texttt{forall x in S: expr} & Set & \texttt{S.map(expr)} & \texttt{forall a in attrs: a.name} \\
\texttt{forall x in S such that P: expr} & Set & \texttt{S.filter(P).map(expr)} & \texttt{forall a in attrs such that a.isPublic: a.name} \\
\texttt{forall x in S do action} & -- & Imperative loop & \texttt{forall a in attrs do createColumn(a)} \\
\bottomrule
\end{tabular}
\end{center}

\subsubsection{Separator Design}

The three separators are \emph{not} interchangeable; each has a precise role:

\begin{itemize}
    \item \kw{such that} introduces a \emph{boolean filter condition}. It reads as a natural qualifier: ``all attributes \emph{such that} they are public.''
    \item \texttt{:} introduces a \emph{value projection}. It reads as ``give me'': ``all attributes: their name'' (i.e., the names of all attributes).
    \item \kw{do} introduces an \emph{imperative action}. It is reserved for \jjtl{} transformation actions and is not available in pure \jjel{} expressions. This separation ensures \jjel{} remains side-effect-free.
\end{itemize}

The decision to use \kw{such that} (two words) rather than a symbol like \texttt{|} was deliberate: it prioritizes readability and matches the set-builder notation commonly taught in discrete mathematics courses ($\{x \in S \mid P(x)\}$ is read ``$x$ in $S$ \emph{such that} $P(x)$'').

\subsubsection{Why Not \texttt{filter} and \texttt{map}?}

The \texttt{filter}/\texttt{map} pair requires the user to:
\begin{enumerate}
    \item Know two distinct method names and their order (filter first, then map).
    \item Write a lambda for each: \texttt{attrs.filter(a => a.isPublic).map(a => a.name)}.
    \item Understand method chaining and intermediate collections.
\end{enumerate}

With \kw{forall}, the same expression becomes:
\begin{lstlisting}[style=jjel,numbers=none]
forall a in attrs such that a.isPublic: a.name
\end{lstlisting}

This is one concept instead of three, with a single binding variable (\texttt{a}) instead of two lambdas. The cognitive load reduction is significant, especially for domain experts who are not programmers.


\subsection{Existential Check: \texttt{exists}}\label{sec:exists}

While \kw{forall} returns a set, \kw{exists} returns a \textbf{boolean}: does at least one element satisfy the condition?

\begin{lstlisting}[style=jjel,numbers=none]
exists a in attributes: a.type == "String"
exists a in attributes such that a.isPublic and not a.isDerived
\end{lstlisting}

\kw{exists} is semantically derived from \kw{forall}:
\[
\texttt{exists x in S: P} \equiv \texttt{(forall x in S such that P).isNotEmpty}
\]

For \kw{exists}, the separators \texttt{:} and \kw{such that} are \emph{synonyms}: both introduce the boolean predicate. This is because \kw{exists} has no projection: it always returns boolean. The synonym design avoids forcing users to remember which separator to use for \kw{exists}.


\subsection{Logical Implication: \texttt{implies}}

\begin{lstlisting}[style=jjel,numbers=none]
isAbstract implies subClasses.isNotEmpty
a.isPublic implies a.type != null
\end{lstlisting}

\kw{implies} is semantically equivalent to \texttt{not P or Q} but reads as natural language. It is right-associative, short-circuit evaluated, and has lower precedence than \kw{and}/\kw{or}.

The motivation is direct: model validation invariants frequently express implications (``if a class is abstract, then it must have subclasses''). Without \kw{implies}, users must negate the antecedent manually (\texttt{not isAbstract or subClasses.isNotEmpty}), which is less readable and more error-prone.


\subsection{Context Binding: \texttt{with...do}}

\begin{lstlisting}[style=jjel,numbers=none]
with selectedNode do name.pascalCase()
with selectedNode do forall a in attributes such that a.isPublic: a.name
\end{lstlisting}

The \kw{with} construct establishes a context object whose properties become directly accessible in the body. This is essential in the Console REPL, where there is no implicit source element (unlike \jjtl{} rules where the matched element provides implicit context).

The choice of \kw{do} as the body separator (rather than a block or colon) was deliberate: it reads naturally (``with this node, \emph{do} [evaluate] this expression'') and parallels the imperative \kw{do} in \kw{forall} by indicating ``what to compute in this context.''


\subsection{Conditional: \texttt{if...then...else}}

\begin{lstlisting}[style=jjel,numbers=none]
if isAbstract then "abstract" else "concrete"
if lowerBound == 0 then "NULL" else "NOT NULL"
\end{lstlisting}

This construct is retained from common expression languages. An alternative syntax (\texttt{when...then...otherwise}) was considered for its more declarative flavor, but \texttt{if/then/else} was kept for its universality and zero learning curve.

When \kw{else} is omitted, the expression returns \texttt{null} for the false branch, maintaining the invariant that every expression produces a value.


\subsection{Lambda Expressions}

\begin{lstlisting}[style=jjel,numbers=none]
a => a.name
(a, b) => a + b
\end{lstlisting}

Lambdas use the \texttt{=>} arrow syntax. An earlier design used \texttt{:} as the lambda separator (\texttt{a: a.name}), but this conflicted with the projection role of \texttt{:} in \kw{forall}. The switch to \texttt{=>} freed \texttt{:} for its set-theoretic role and made the grammar LL(1)-parseable without backtracking.

Lambdas are primarily used with collection methods not subsumed by \kw{forall}: \texttt{sortBy}, \texttt{groupBy}, \texttt{distinctBy}, etc.


\subsection{Null Safety}

\begin{lstlisting}[style=jjel,numbers=none]
parent?.name                     -- null if parent is null
parent?.name ?? "no parent"      -- "no parent" if chain is null
element?.container?.package?.name ?? "default"
\end{lstlisting}

Two operators handle null propagation:
\begin{itemize}
    \item \texttt{?.} (null-safe navigation): returns \texttt{null} if the receiver is null, without throwing.
    \item \texttt{??} (null coalesce): returns the right operand if the left is null.
\end{itemize}

Using \texttt{.} on a null value throws an error by design: this forces users to explicitly acknowledge potential nulls via \texttt{?.}, making null handling visible rather than silently propagating.


\subsection{Operator Precedence}

\begin{center}
\begin{tabular}{@{}clc@{}}
\toprule
\textbf{Level} & \textbf{Construct} & \textbf{Associativity} \\
\midrule
1 (lowest) & \texttt{if ... then ... else ...} & right \\
2 & \texttt{??} (null coalesce) & left \\
3 & \texttt{implies} & right \\
4 & \texttt{or} & left \\
5 & \texttt{and} & left \\
6 & \texttt{==}, \texttt{!=} & left \\
7 & \texttt{<}, \texttt{>}, \texttt{<=}, \texttt{>=} & left \\
8 & \texttt{is} & -- \\
9 & \texttt{+}, \texttt{-} & left \\
10 & \texttt{*}, \texttt{/}, \texttt{\%} & left \\
11 & \texttt{not}, \texttt{-} (unary) & right \\
12 (highest) & \texttt{.}, \texttt{?.}, \texttt{[index]} & left \\
\bottomrule
\end{tabular}
\end{center}


\subsection{Built-in Methods}

\jjel{} provides over 125 built-in methods organized by type. The guiding principle is that the most frequent MDE operations should be native, requiring no import or boilerplate.

\begin{center}
\begin{tabular}{@{}lrl@{}}
\toprule
\textbf{Category} & \textbf{Count} & \textbf{Highlights} \\
\midrule
String & 33 & \texttt{camelCase()}, \texttt{pascalCase()}, \texttt{snakeCase()}, \texttt{kebabCase()} \\
Number & 30 & \texttt{abs()}, \texttt{clamp()}, \texttt{between()}, trigonometry \\
Collection & 28 & \texttt{sortBy()}, \texttt{groupBy()}, \texttt{flatten()}, \texttt{distinct()} \\
Date & 39 & \texttt{addDays()}, \texttt{diffMonths()}, \texttt{format()} \\
\bottomrule
\end{tabular}
\end{center}

Naming convention methods (\texttt{camelCase}, \texttt{pascalCase}, \texttt{snakeCase}, \texttt{kebabCase}) are native because they are essential in model transformations (e.g., converting class names to table names, attribute names to column names). No other transformation language provides these natively.

Note: \texttt{filter()} and \texttt{map()} are \emph{removed} from the collection methods, as they are replaced by the \kw{forall} construct. The remaining collection methods are those that \kw{forall} does not subsume: aggregation (\texttt{sum}, \texttt{avg}), ordering (\texttt{sortBy}), structure (\texttt{groupBy}, \texttt{flatten}), and boolean checks (\texttt{all}, \texttt{any}, \texttt{none}).


\subsection{Evaluation Rules}

\subsubsection{Short-Circuit Evaluation}

\kw{and} and \kw{or} use \textbf{short-circuit} evaluation: the right operand is not evaluated if the result is determined by the left operand. This is essential for guard-like expressions:

\begin{lstlisting}[style=jjel,numbers=none]
parent != null and parent.name.isNotEmpty    -- safe: parent.name not reached if null
\end{lstlisting}

\subsubsection{Truthiness}

\jjel{} uses a broad truthiness model: \texttt{null}, \texttt{false}, \texttt{0}, \texttt{""} (empty string), and \texttt{[]} (empty array) are all falsy. This enables concise guards (\texttt{if attributes then ...} instead of \texttt{if attributes.isNotEmpty then ...}).

\subsubsection{Error Handling}

\jjel{} distinguishes between recoverable and non-recoverable errors:
\begin{itemize}
    \item \textbf{Silent null}: undefined variables, property not found on object (with console warning and typo suggestions via Levenshtein distance), division by zero, type mismatches in arithmetic.
    \item \textbf{Thrown error}: \texttt{.property} on \texttt{null} (forces use of \texttt{?.}), unknown method call (with suggestions).
\end{itemize}


% ============================================
% 3. JjTL
% ============================================
\section{JjTL: Jjodel Transformation Language}\label{sec:jjtl}

\jjtl{} is a declarative, rule-based Model-to-Model transformation language. It uses \jjel{} as its expression sub-language.

\subsection{Design Philosophy}

\jjtl{} targets the ``sweet spot'' between ATL's complexity and manual Java code:

\begin{itemize}
    \item \textbf{Simpler than ATL}: No two-phase execution, no \texttt{thisModule} reference, no distinction between matched/lazy/called rules. Every rule is a class mapping.
    \item \textbf{More structured than ETL}: Attribute mappings use a declarative binding syntax (\texttt{->}) rather than imperative assignment. Guards are explicit (\kw{when}) rather than statement blocks.
    \item \textbf{JjEL instead of OCL}: Expressions use \jjel{}'s set-theoretic semantics, null-safe navigation, and native naming conventions, features that OCL does not provide.
\end{itemize}

\subsection{Transformation Structure}

A \jjtl{} transformation declares its name, source metamodel, target metamodel, and a set of class mappings and helper functions:

\begin{lstlisting}[style=jjtl]
transformation StateMachine2PetriNet

from StateMachineMM
to   PetriNetMM

# Class mappings follow
State -> Place {
    name -> label
    isInitial -> tokens : true=1, false=0
}

# Helper functions
helper formatLabel(s: String) -> String {
    s.toUpper()
}
\end{lstlisting}


\subsection{Class Mappings}

A class mapping specifies a correspondence between a source class and a target class:

\begin{lstlisting}[style=jjtl,numbers=none]
SourceClass -> TargetClass [multiplicity] when { guard } {
    attributeMappings...
}
\end{lstlisting}

\subsubsection{Multiplicity}

Multiplicity controls how many target instances are created per source instance:
\begin{itemize}
    \item No multiplicity: one-to-one (default).
    \item \texttt{[*]}: unbounded (used with collection sources).
    \item \texttt{[n]}: exactly $n$ target instances per source.
    \item \texttt{[m..n]}: between $m$ and $n$ instances.
\end{itemize}

\subsubsection{Guard Conditions}

The \kw{when} clause filters which source instances are transformed. The guard is a \jjel{} expression evaluated with the source instance as implicit context:

\begin{lstlisting}[style=jjtl,numbers=none]
Class -> Table when { not isAbstract and attributes.isNotEmpty } {
    ...
}
\end{lstlisting}

The implicit context means property access is unqualified: \texttt{isAbstract} refers to \texttt{source.isAbstract}. This reduces verbosity compared to ATL (where the guard must use the bound variable: \texttt{c.isAbstract}) and QVT-O (where \texttt{self.isAbstract} is required).


\subsection{Attribute Mappings}

\subsubsection{Direct Mapping}

The simplest form copies a value from source to target:
\begin{lstlisting}[style=jjtl,numbers=none]
name -> tableName
\end{lstlisting}

\subsubsection{Value Mapping (Enum Conversion)}

Discrete value conversions use the \texttt{:} separator followed by comma-separated pairs:
\begin{lstlisting}[style=jjtl,numbers=none]
isInitial -> tokens : true=1, false=0
status -> code : "active"=1, "inactive"=0, "deleted"=-1
\end{lstlisting}

Unmapped values pass through unchanged (fallback identity).

\subsubsection{Expression Mapping}

Arbitrary \jjel{} expressions compute the target value:
\begin{lstlisting}[style=jjtl,numbers=none]
name -> tableName : name.snakeCase()
-> columnDefs : forall a in attributes: a.name + " " + a.type.toUpper()
\end{lstlisting}

\subsubsection{Object Creation}

Nested target objects are created inline:
\begin{lstlisting}[style=jjtl,numbers=none]
-> inputArcs {
    -> Arc {
        place -> source
        weight -> 1
    }
}
\end{lstlisting}


\subsection{Helper Functions}

Helpers are named, reusable \jjel{} expressions:
\begin{lstlisting}[style=jjtl,numbers=none]
helper formatLabel(name: String, prefix: String) -> String {
    prefix + "_" + name.snakeCase()
}
\end{lstlisting}

Helpers are registered in the evaluation context and callable from any expression.


\subsection{Trace Model}

\jjtl{} automatically records a trace model during execution, linking each source element to its target element(s) with binding-level detail:

\begin{itemize}
    \item \textbf{Rule-level trace}: which class mapping produced which target instance.
    \item \textbf{Binding-level trace}: for each attribute mapping, the source value, target value, expression used, and whether the mapping is \emph{invertible}.
\end{itemize}

\subsubsection{Invertibility Analysis}

\jjtl{} automatically determines whether each attribute mapping can be reversed:
\begin{itemize}
    \item Direct mappings (\texttt{a -> b}): always invertible.
    \item Value mappings (\texttt{true=1, false=0}): invertible if and only if the mapping is injective (no duplicate target values).
    \item Expression mappings: invertible only for simple property references; function calls and arithmetic are marked non-invertible.
\end{itemize}


\subsection{Interactive Features}

\jjtl{} supports user interaction during transformation execution via four built-in statements:

\begin{center}
\begin{tabular}{@{}lll@{}}
\toprule
\textbf{Construct} & \textbf{Behavior} & \textbf{Returns} \\
\midrule
\texttt{alert(msg, type)} & Blocking modal & -- \\
\texttt{notify(msg, duration)} & Non-blocking toast & -- \\
\texttt{prompt(msg, default)} & Text input dialog & String \\
\texttt{input(msg, type, options)} & Typed input dialog & Varies \\
\bottomrule
\end{tabular}
\end{center}

These enable semi-automated transformations where certain decisions require human input, a feature absent from ATL, ETL, QVT-R, and QVT-O.


\subsection{Execution Model}

The \jjtl{} executor follows a simple pipeline:

\begin{enumerate}
    \item \textbf{Initialize}: bind source model elements in the evaluation context; register helpers.
    \item \textbf{Match}: for each class mapping, find all source instances of the matching type.
    \item \textbf{Guard}: filter instances through the \kw{when} condition.
    \item \textbf{Create}: instantiate target elements according to multiplicity.
    \item \textbf{Bind}: evaluate attribute mappings and assign values.
    \item \textbf{Trace}: record the trace link with invertibility information.
\end{enumerate}

Unlike ATL's two-phase execution (where all bindings are computed before any imperative blocks run), \jjtl{} processes each rule atomically: match, create, bind. This simplifies reasoning about execution order.


% ============================================
% 4. COMPARISON
% ============================================
\section{Comparative Analysis}\label{sec:comparison}

We now compare \jjtl{}/\jjel{} with four established transformation languages, examining both syntactic design and underlying philosophy.

\subsection{ATL (Atlas Transformation Language)}

\subsubsection{Overview}

ATL is a hybrid declarative/imperative M2M transformation language developed at INRIA as part of the ATLAS group. It uses OCL as its expression sub-language and has been the most widely used academic transformation language since the mid-2000s.

\subsubsection{Structural Comparison}

\begin{center}
\begin{tabularx}{\textwidth}{@{}lXX@{}}
\toprule
\textbf{Aspect} & \textbf{ATL} & \textbf{JjTL} \\
\midrule
Header & \texttt{module Name; create Out from In} & \texttt{transformation Name \textbackslash n from X \textbackslash n to Y} \\
Rule type & Matched / Lazy / Called / Entrypoint & Single type: class mapping \\
Binding & \texttt{targetAttr <- expr} & \texttt{sourceAttr -> targetAttr} \\
Guard & OCL expression in \texttt{from} clause & \kw{when} block with \jjel{} \\
Trace resolution & Implicit via matched rules & Implicit (like ATL) with per-binding invertibility \\
Imperative block & \texttt{do \{ ... \}} (post-binding) & \kw{do} in \kw{forall} (per-element) \\
Expression language & OCL & \jjel{} \\
\bottomrule
\end{tabularx}
\end{center}

\subsubsection{Key Design Differences}

\paragraph{Binding direction.} ATL uses \texttt{<-} (``target receives from expression''), while \jjtl{} uses \texttt{->} (``source maps to target''). The \jjtl{} direction mirrors the natural reading of transformation rules: the \emph{source} drives the transformation. ATL's direction mirrors assignment (\texttt{x := expr}).

\paragraph{Rule taxonomy.} ATL distinguishes four rule types (matched, lazy, called, entrypoint), each with different execution semantics. \jjtl{} has a single rule type with optional multiplicity and guard, reducing the conceptual surface.

\paragraph{Two-phase execution.} ATL's matched rules execute in two phases: first, all target elements are created and bindings computed (declarative phase); then, \texttt{do} blocks run (imperative phase). This design enables inter-rule references during binding but makes execution order non-obvious. \jjtl{} processes each mapping atomically: match, create, bind.

\paragraph{Expression language.} ATL uses OCL, which lacks null-safe navigation, naming convention methods, and set comprehension via \kw{forall}. OCL collection operations (\texttt{->select()}, \texttt{->collect()}) require lambda syntax (\texttt{e | expr}) and explicit collection type management (\texttt{Set}, \texttt{Sequence}, \texttt{Bag}, \texttt{OrderedSet}).


\subsection{EOL / ETL (Epsilon Platform)}

\subsubsection{Overview}

EOL is the imperative base language of the Epsilon platform; ETL extends it for M2M transformations. EOL takes a pragmatic approach, designed to be familiar to Java/JavaScript developers.

\subsubsection{Structural Comparison}

\begin{center}
\begin{tabularx}{\textwidth}{@{}lXX@{}}
\toprule
\textbf{Aspect} & \textbf{ETL} & \textbf{JjTL} \\
\midrule
Rule syntax & \texttt{rule Name transform s:T to t:T \{ ... \}} & \texttt{S -> T \{ ... \}} \\
Binding & Imperative assignment: \texttt{t.x = s.y} & Declarative: \texttt{y -> x} \\
Guard & \texttt{guard: expr} or \texttt{guard \{ block \}} & \texttt{when \{ expr \}} \\
Trace resolution & \texttt{::=} operator / \texttt{equivalent()} & Implicit with per-binding invertibility \\
Rule modifiers & \texttt{@abstract}, \texttt{@lazy}, \texttt{@primary}, \texttt{@greedy} & Multiplicity (\texttt{[*]}) \\
Inheritance & \texttt{extends parentRule} & Not supported \\
Control flow & \texttt{for}, \texttt{if/else}, \texttt{while}, \texttt{switch} & \kw{forall}, \texttt{if/then/else} \\
\bottomrule
\end{tabularx}
\end{center}

\subsubsection{Key Design Differences}

\paragraph{Declarative vs.\ imperative bindings.} ETL attribute mappings are imperative assignments (\texttt{t.name = s.name.toUpperCase()}). \jjtl{} attribute mappings are declarative bindings (\texttt{name -> tableName : name.toUpper()}). The declarative approach enables automatic invertibility analysis, which is impossible with imperative assignments.

\paragraph{Verbosity.} ETL requires explicit variable declarations, \texttt{new} for object creation, and method call syntax for collection operations. \jjtl{} infers targets from the rule header, creates objects implicitly, and uses \kw{forall} instead of explicit loops or lambda chains.

\paragraph{Rule inheritance.} ETL supports \texttt{extends} for rule inheritance (abstract rules with shared logic). \jjtl{} does not have rule inheritance, favoring flat, self-contained rules. This is a deliberate simplicity trade-off.


\subsection{QVT Relations (QVT-R)}

\subsubsection{Overview}

QVT-R is the declarative relational sub-language of the OMG QVT standard. Its distinguishing features are bidirectional specification and pattern-based relation declarations.

\subsubsection{Structural Comparison}

\begin{center}
\begin{tabularx}{\textwidth}{@{}lXX@{}}
\toprule
\textbf{Aspect} & \textbf{QVT-R} & \textbf{JjTL} \\
\midrule
Paradigm & Fully declarative, bidirectional & Declarative, bidirectional (with unidirectional extension) \\
Rule structure & \texttt{relation} with domain patterns & Class mapping with attribute bindings \\
Directionality & \texttt{checkonly} / \texttt{enforce} per domain & Bidirectional core; unidirectional extension \\
Precondition & \texttt{when \{ ... \}} (relation calls) & \texttt{when \{ JjEL expr \}} \\
Postcondition & \texttt{where \{ ... \}} (relation calls) & Not explicit (sequential) \\
Expression language & OCL & \jjel{} \\
Variable binding & Implicit via pattern variables & Explicit via source property names \\
\bottomrule
\end{tabularx}
\end{center}

\subsubsection{Key Design Differences}

\paragraph{Bidirectionality.} QVT-R's core promise is that the same transformation can be executed in either direction. In practice, this creates significant complexity: users must reason about enforcement behavior in both directions simultaneously, and non-bijective mappings (common in real transformations like class-to-table) introduce semantic ambiguities. \jjtl{}'s declarative core is also bidirectional, but takes a different approach: invertibility is tracked per-binding rather than per-transformation, giving fine-grained control over which parts of a mapping can be reversed. Additionally, \jjtl{} provides a unidirectional extension for cases where source-to-target-only semantics are needed.

\paragraph{Pattern matching vs.\ explicit mapping.} QVT-R uses pattern-based domain declarations where variables are bound across domains by name equality. This is elegant but hard to read: understanding which variables bind to what requires scanning the entire relation. \jjtl{} uses explicit \texttt{->} mappings, making the data flow immediately visible.

\paragraph{Relation calls.} QVT-R's \texttt{when}/\texttt{where} clauses invoke other relations, creating a web of inter-dependent correspondences. This compositional power is also a source of complexity: debugging requires tracing through chains of relation calls. \jjtl{} guards are self-contained \jjel{} expressions.


\subsection{QVT Operational (QVT-O)}

\subsubsection{Overview}

QVT-O is the imperative sub-language of the QVT standard. It uses \texttt{mapping} operations with automatic trace management, OCL-based expressions, and a three-section body (\texttt{init}/population/\texttt{end}).

\subsubsection{Structural Comparison}

\begin{center}
\begin{tabularx}{\textwidth}{@{}lXX@{}}
\toprule
\textbf{Aspect} & \textbf{QVT-O} & \textbf{JjTL} \\
\midrule
Entry point & Explicit \texttt{main()} & Implicit (all rules fire) \\
Mapping syntax & \texttt{mapping T::op() : TT} & \texttt{S -> T \{ ... \}} \\
Implicit context & \texttt{self} & Unqualified (source properties) \\
Object creation & \texttt{object Type \{ ... \}} & Nested \texttt{-> Type \{ ... \}} \\
Guard & \texttt{when \{ ... \}} & \texttt{when \{ ... \}} \\
Trace resolution & Implicit via mapping invocation & Implicit with per-binding invertibility \\
Assignment & \texttt{:=} & \texttt{->} \\
Expression language & OCL & \jjel{} \\
\bottomrule
\end{tabularx}
\end{center}

\subsubsection{Key Design Differences}

\paragraph{Entry point.} QVT-O requires an explicit \texttt{main()} that orchestrates the transformation by navigating the source model and invoking mappings. \jjtl{} fires all class mappings automatically for matching source instances (pattern-matched, like ATL). This reduces boilerplate but limits control over execution order.

\paragraph{Self vs.\ implicit context.} QVT-O uses \texttt{self} to refer to the source element within a mapping. \jjtl{} uses an implicit context: properties of the source element are accessed without qualification. This is a significant verbosity reduction but requires clear scoping rules.


\subsection{Feature Comparison Matrix}

\begin{center}
\small
\begin{tabularx}{\textwidth}{@{}l*{5}{c}@{}}
\toprule
\textbf{Feature} & \textbf{ATL} & \textbf{ETL} & \textbf{QVT-R} & \textbf{QVT-O} & \textbf{JjTL} \\
\midrule
Declarative rules & \checkmark & -- & \checkmark & -- & \checkmark \\
Imperative blocks & \checkmark & \checkmark & -- & \checkmark & Limited$^a$ \\
Bidirectional & -- & -- & \checkmark & -- & \checkmark$^e$ \\
Automatic trace & \checkmark & \checkmark & \checkmark & \checkmark & \checkmark \\
Invertibility tracking & -- & -- & -- & -- & \checkmark \\
Guard conditions & \checkmark & \checkmark & \checkmark$^b$ & \checkmark & \checkmark \\
Multiplicity & -- & -- & -- & -- & \checkmark \\
Null-safe navigation & -- & \checkmark$^c$ & -- & -- & \checkmark \\
Naming conventions & -- & -- & -- & -- & \checkmark \\
Interactive (user input) & -- & -- & -- & -- & \checkmark \\
Set comprehension (\kw{forall}) & -- & -- & -- & -- & \checkmark \\
Rule inheritance & -- & \checkmark & -- & \checkmark$^d$ & -- \\
Expression language & OCL & EOL & OCL & OCL & \jjel{} \\
\bottomrule
\end{tabularx}

\vspace{4pt}
\footnotesize
$^a$ Via \kw{forall...do} in expressions. $^b$ Via \texttt{when}/\texttt{where} clauses. $^c$ Since Epsilon 2.1. $^d$ Via mapping inheritance. $^e$ Declarative core is bidirectional; also has a unidirectional extension.
\end{center}


\subsection{Expression Language Comparison}

A critical differentiator is the expression sub-language. We compare equivalent expressions across OCL (used by ATL, QVT-R, QVT-O), EOL (used by ETL), and \jjel{} (used by \jjtl{}):

\begin{center}
\small
\begin{tabularx}{\textwidth}{@{}lXXX@{}}
\toprule
\textbf{Task} & \textbf{OCL} & \textbf{EOL} & \textbf{JjEL} \\
\midrule
Public attr. names &
\texttt{self.attr->select(a|a.isPublic)->collect(a|a.name)} &
\texttt{self.attr.select(a|a.isPublic).collect(a|a.name)} &
\texttt{forall a in attr such that a.isPublic: a.name} \\
\addlinespace
Exists string attr &
\texttt{self.attr->exists(a|a.type.name='String')} &
\texttt{self.attr.exists(a|a.type.name='String')} &
\texttt{exists a in attr: a.type == "String"} \\
\addlinespace
Abstract $\Rightarrow$ has children &
\texttt{self.isAbstract implies self.subClasses->notEmpty()} &
\texttt{self.isAbstract implies self.subClasses.notEmpty()} &
\texttt{isAbstract implies subClasses.isNotEmpty} \\
\addlinespace
Null-safe name &
\texttt{if self.parent.oclIsUndefined() then 'none' else self.parent.name endif} &
\texttt{self.parent?.name ?: 'none'} &
\texttt{parent?.name ?? "none"} \\
\addlinespace
To snake\_case &
(no native support) &
(no native support) &
\texttt{name.snakeCase()} \\
\bottomrule
\end{tabularx}
\end{center}

Key observations:
\begin{enumerate}
    \item \jjel{}'s \kw{forall} replaces the \texttt{->select()->collect()} chain with a single construct.
    \item OCL's null handling requires verbose \texttt{oclIsUndefined()} checks.
    \item Naming convention methods (\texttt{snakeCase}, \texttt{camelCase}) are unique to \jjel{}.
    \item \jjel{} uses implicit context (no \texttt{self}), reducing token count.
\end{enumerate}


% ============================================
% 5. EXAMPLES
% ============================================
\section{Transformation Examples}\label{sec:examples}

We present three transformations in parallel across ATL, ETL, and \jjtl{}, each demonstrating different language features.


\subsection{Example 1: UML Class to Relational Table}

This classic transformation maps UML classes to relational tables: each non-abstract class becomes a table, each single-valued attribute becomes a column, and each multi-valued attribute becomes a junction table.

\subsubsection{ATL}

\begin{lstlisting}[style=atl,caption={Class2Relational in ATL}]
module Class2Relational;
create OUT : Relational from IN : Class;

helper context Class!Attribute def: isSingleValued() : Boolean =
    not self.multiValued;

rule Class2Table {
    from
        c : Class!Class (
            not c.isAbstract
        )
    to
        t : Relational!Table (
            name <- c.name,
            col  <- c.attr->select(a | a.isSingleValued())
                           ->collect(a | thisModule.Attr2Col(a))
                           ->union(Sequence{key}),
            key  <- Set{key}
        ),
        key : Relational!Column (
            name <- 'Id'
        )
}

lazy rule Attr2Col {
    from
        a : Class!Attribute
    to
        c : Relational!Column (
            name <- a.name,
            type <- a.type.name
        )
}

rule MultiAttr2Table {
    from
        a : Class!Attribute (
            a.multiValued
        )
    to
        t : Relational!Table (
            name <- a.owner.name + '_' + a.name,
            col  <- Sequence{id, val}
        ),
        id : Relational!Column (
            name <- 'Id'
        ),
        val : Relational!Column (
            name <- a.name,
            type <- a.type.name
        )
}
\end{lstlisting}

\subsubsection{ETL}

\begin{lstlisting}[style=etl,caption={Class2Relational in ETL}]
rule Class2Table
    transform c : Class!Class
    to t : Relational!Table {

    guard : not c.isAbstract

    t.name = c.name;

    // Primary key
    var pk = new Relational!Column;
    pk.name = "Id";
    pk.type = "INTEGER";
    t.col.add(pk);

    // Single-valued attributes
    for (a in c.attr.select(a | not a.multiValued)) {
        var col = new Relational!Column;
        col.name = a.name;
        col.type = a.type.name;
        t.col.add(col);
    }
}

rule MultiAttr2Table
    transform a : Class!Attribute
    to t : Relational!Table {

    guard : a.multiValued

    t.name = a.owner.name + "_" + a.name;

    var id = new Relational!Column;
    id.name = "Id";
    t.col.add(id);

    var val = new Relational!Column;
    val.name = a.name;
    val.type = a.type.name;
    t.col.add(val);
}
\end{lstlisting}

\subsubsection{JjTL}

\begin{lstlisting}[style=jjtl,caption={Class2Relational in JjTL}]
transformation Class2Relational

from ClassMM
to   RelationalMM

# Non-abstract classes become tables
Class -> Table when { not isAbstract } {
    name -> name : name.snakeCase()

    # Primary key column
    -> col {
        -> Column {
            -> name : "id"
            -> type : "INTEGER"
        }
    }

    # Single-valued attributes become columns
    -> col {
        forall a in attributes such that not a.multiValued
        -> Column {
            a.name -> name : a.name.snakeCase()
            a.type.name -> type
        }
    }
}

# Multi-valued attributes become junction tables
Attribute -> Table when { multiValued } {
    -> name : owner.name.snakeCase() + "_" + name.snakeCase()

    -> col {
        -> Column {
            -> name : "id"
        }
    }
    -> col {
        -> Column {
            name -> name : name.snakeCase()
            type.name -> type
        }
    }
}
\end{lstlisting}


\subsection{Example 2: State Machine to Petri Net}

This transformation converts a state machine (states and transitions) into a Petri net (places, transitions, and arcs).

\subsubsection{ATL}

\begin{lstlisting}[style=atl,caption={StateMachine2PetriNet in ATL}]
module SM2PN;
create OUT : PetriNet from IN : StateMachine;

rule State2Place {
    from
        s : StateMachine!State
    to
        p : PetriNet!Place (
            name   <- s.name,
            tokens <- if s.isInitial then 1 else 0 endif
        )
}

rule Transition2Transition {
    from
        t : StateMachine!Transition
    to
        pt : PetriNet!Transition (
            name <- t.label
        ),
        inArc : PetriNet!Arc (
            source    <- t.source,  -- resolved via trace
            target    <- pt,
            weight    <- 1
        ),
        outArc : PetriNet!Arc (
            source    <- pt,
            target    <- t.target,  -- resolved via trace
            weight    <- 1
        )
}
\end{lstlisting}

\subsubsection{ETL}

\begin{lstlisting}[style=etl,caption={StateMachine2PetriNet in ETL}]
rule State2Place
    transform s : SM!State
    to p : PN!Place {

    p.name   = s.name;
    p.tokens = if (s.isInitial) 1 else 0;
}

rule Transition2Transition
    transform t : SM!Transition
    to pt : PN!Transition {

    pt.name = t.label;

    // Input arc: from source place to this transition
    var inArc = new PN!Arc;
    inArc.source = t.source.equivalent();
    inArc.target = pt;
    inArc.weight = 1;

    // Output arc: from this transition to target place
    var outArc = new PN!Arc;
    outArc.source = pt;
    outArc.target = t.target.equivalent();
    outArc.weight = 1;
}
\end{lstlisting}

\subsubsection{JjTL}

\begin{lstlisting}[style=jjtl,caption={StateMachine2PetriNet in JjTL}]
transformation StateMachine2PetriNet

from StateMachineMM
to   PetriNetMM

State -> Place {
    name -> name
    isInitial -> tokens : true=1, false=0
}

Transition -> Transition [*] {
    label -> name

    # Input arc: source state -> this transition
    -> inputArcs {
        -> Arc {
            source -> source
            -> weight : 1
        }
    }

    # Output arc: this transition -> target state
    -> outputArcs {
        -> Arc {
            target -> target
            -> weight : 1
        }
    }
}
\end{lstlisting}

\paragraph{Comparison notes.}
\begin{itemize}
    \item ATL and \jjtl{} both use implicit trace resolution: when a source element (e.g., a \texttt{State}) is assigned to a target property that expects the corresponding target type (e.g., a \texttt{Place}), the engine automatically resolves the trace link. This keeps the mapping syntax clean and declarative.
    \item ETL requires explicit \texttt{.equivalent()} calls to resolve trace links.
    \item The value mapping \texttt{true=1, false=0} in \jjtl{} replaces ATL's verbose \texttt{if...then...else...endif}.
\end{itemize}


\subsection{Example 3: UML to Java (Structural)}

This example maps UML classes to Java class definitions, demonstrating naming conventions, inheritance, and collection processing.

\subsubsection{ATL}

\begin{lstlisting}[style=atl,caption={UML2Java structural mapping in ATL}]
module UML2Java;
create OUT : Java from IN : UML;

helper context UML!Class def: javaName() : String =
    self.name;
    -- Note: no native camelCase/PascalCase in OCL

helper context UML!Attribute def: getterName() : String =
    'get' + self.name;  -- cannot capitalize first letter natively

rule Class2JavaClass {
    from
        c : UML!Class
    to
        j : Java!JavaClass (
            name       <- c.javaName(),
            isAbstract <- c.isAbstract,
            superClass <- c.superClass,  -- trace-resolved
            fields     <- c.attributes->select(a |
                            not a.isDerived),
            methods    <- c.attributes->select(a |
                            not a.isDerived)
                            ->collect(a | thisModule.Attr2Getter(a))
        )
}

lazy rule Attr2Getter {
    from
        a : UML!Attribute
    to
        m : Java!Method (
            name       <- a.getterName(),
            returnType <- a.type,
            body       <- 'return this.' + a.name + ';'
        )
}
\end{lstlisting}

\subsubsection{ETL}

\begin{lstlisting}[style=etl,caption={UML2Java structural mapping in ETL}]
rule Class2JavaClass
    transform c : UML!Class
    to j : Java!JavaClass {

    j.name       = c.name;
    j.isAbstract = c.isAbstract;
    j.superClass ::= c.superClass;

    for (a in c.attributes.select(a | not a.isDerived)) {
        var field = new Java!Field;
        field.name = a.name;
        field.type ::= a.type;
        j.fields.add(field);

        var getter = new Java!Method;
        getter.name = "get" + a.name.firstToUpperCase();
        getter.returnType ::= a.type;
        getter.body = "return this." + a.name + ";";
        j.methods.add(getter);
    }
}
\end{lstlisting}

\subsubsection{JjTL}

\begin{lstlisting}[style=jjtl,caption={UML2Java structural mapping in JjTL}]
transformation UML2Java

from UML_MM
to   JavaMM

Class -> JavaClass {
    name -> name : name.pascalCase()
    isAbstract -> isAbstract
    superClass -> superClass

    # Fields from non-derived attributes
    -> fields {
        forall a in attributes such that not a.isDerived
        -> Field {
            a.name -> name : a.name.camelCase()
            a.type -> type
        }
    }

    # Getter methods
    -> methods {
        forall a in attributes such that not a.isDerived
        -> Method {
            -> name : "get" + a.name.pascalCase()
            a.type -> returnType
            -> body : "return this." + a.name.camelCase() + ";"
        }
    }
}
\end{lstlisting}

\paragraph{Comparison notes.}
\begin{itemize}
    \item \jjtl{}'s native \texttt{pascalCase()} and \texttt{camelCase()} methods eliminate the need for manual string manipulation helpers that ATL and ETL require.
    \item ATL needs a lazy rule (\texttt{Attr2Getter}) invoked via \texttt{thisModule.Attr2Getter(a)}; \jjtl{} uses inline object creation.
    \item ETL's imperative loop with explicit \texttt{new} and \texttt{.add()} is more verbose than \jjtl{}'s declarative \kw{forall} with nested object creation.
    \item OCL's string handling is the weakest: there is no native method to capitalize the first letter, convert to camelCase, etc.
\end{itemize}


\subsection{Example 4: Validation Invariants}

This example demonstrates \jjel{}'s expressiveness for model validation, compared to OCL.

\begin{center}
\small
\begin{tabularx}{\textwidth}{@{}p{3.5cm}XX@{}}
\toprule
\textbf{Invariant} & \textbf{OCL} & \textbf{JjEL} \\
\midrule
All classes have $\geq$1 attribute &
\texttt{Class.allInstances()->forAll(c | c.attributes->notEmpty())} &
\texttt{forall c in classes: c.attributes.isNotEmpty} \\
\addlinespace
Abstract $\Rightarrow$ has subclasses &
\texttt{self.isAbstract implies self.subClasses->notEmpty()} &
\texttt{isAbstract implies subClasses.isNotEmpty} \\
\addlinespace
Unique class names &
\texttt{Class.allInstances()->isUnique(c | c.name)} &
\texttt{(forall c in classes: c.name).distinct().size == classes.size} \\
\addlinespace
No attribute with null type &
\texttt{self.attributes->forAll(a | not a.type.oclIsUndefined())} &
\texttt{forall a in attributes: a.type != null} \\
\addlinespace
Public attributes $\Rightarrow$ typed &
\texttt{self.attributes->select(a | a.isPublic)->forAll(a | not a.type.oclIsUndefined())} &
\texttt{forall a in attributes such that a.isPublic: a.type != null} \\
\bottomrule
\end{tabularx}
\end{center}

In \jjel{}, the same \kw{forall} construct serves for both selection and projection, while OCL requires different operations (\texttt{->forAll()} for assertion, \texttt{->select()} for filtering, \texttt{->collect()} for mapping). Note also that in \jjel{} validation, \kw{forall} with \texttt{:} returns a set of boolean values; the invariant holds if all values are \texttt{true}. An explicit \texttt{.all(x => x)} or wrapping with \kw{forall ... such that} can make the boolean intent more explicit.


\subsection{Example 5: Guard Conditions}

Guards in transformation rules demonstrate how expression languages handle complex selection criteria.

\begin{center}
\small
\begin{tabularx}{\textwidth}{@{}p{2.8cm}XXX@{}}
\toprule
\textbf{Guard} & \textbf{ATL (OCL)} & \textbf{ETL (EOL)} & \textbf{JjTL (JjEL)} \\
\midrule
Non-abstract with attributes &
\texttt{not c.isAbstract and c.attr->notEmpty()} &
\texttt{not c.isAbstract and c.attr.notEmpty()} &
\texttt{not isAbstract and attributes.isNotEmpty} \\
\addlinespace
Has a string attribute &
\texttt{c.attr->exists(a | a.type.name = 'String')} &
\texttt{c.attr.exists(a | a.type.name = 'String')} &
\texttt{exists a in attributes: a.type == "String"} \\
\addlinespace
All references are navigable &
\texttt{c.refs->forAll(r | r.isNavigable)} &
\texttt{c.refs.forAll(r | r.isNavigable)} &
\texttt{forall r in references: r.isNavigable} \\
\bottomrule
\end{tabularx}
\end{center}


% ============================================
% 6. DISCUSSION
% ============================================
\section{Discussion}\label{sec:discussion}

\subsection{Strengths of the JjTL/JjEL Approach}

\begin{enumerate}
    \item \textbf{Conceptual unification via \kw{forall}}: By giving \kw{forall} set-theoretic semantics, \jjel{} replaces three distinct concepts (filter, map, logical quantification) with one. This is a genuine reduction in cognitive load, not merely syntactic sugar.

    \item \textbf{Domain-native operations}: Built-in naming convention methods (\texttt{camelCase}, \texttt{snakeCase}, etc.) and null-safe navigation address the most common pain points in model transformations, which no other transformation language handles natively.

    \item \textbf{Invertibility tracking}: \jjtl{}'s automatic classification of attribute mappings as invertible or non-invertible is unique among the compared languages and enables downstream applications like bidirectional synchronization and transformation debugging.

    \item \textbf{Interactive transformations}: The \texttt{prompt}/\texttt{input} mechanism supports semi-automated scenarios that other languages cannot express.

    \item \textbf{Single execution model}: Unlike ATL's two-phase matched rules or QVT-R's complex enforcement semantics, \jjtl{}'s execution model is straightforward: match, create, bind.
\end{enumerate}

\subsection{Limitations and Trade-offs}

\begin{enumerate}
    \item \textbf{Limited imperative support}: \jjtl{}'s \texttt{do} blocks provide basic imperative capability, but lack the full expressiveness of ATL's called rules or QVT-O's mapping operations for complex procedural logic.

    \item \textbf{No rule inheritance}: Unlike ETL's \texttt{extends} and QVT-O's mapping inheritance, \jjtl{} rules are flat. Complex transformations with shared logic across rules must use helper functions.

    \item \textbf{No explicit entry point}: The automatic ``fire all matching rules'' execution model (like ATL) limits control over rule ordering. QVT-O's explicit \texttt{main()} provides more control.

    \item \textbf{Maturity}: ATL and ETL have decades of academic use, extensive example repositories, and battle-tested implementations. \jjtl{}/\jjel{} are new and will require the same maturation process.
\end{enumerate}

\subsection{Open Design Points}

Several aspects of \jjel{} are under active review:

\begin{enumerate}
    \item \textbf{\texttt{if/then/else} syntax}: The conditional has a procedural flavor. Alternatives like \texttt{when/then/otherwise} have been considered for their more declarative reading.

    \item \textbf{String interpolation}: The lexer supports \texttt{\$\{...\}} syntax but the parser does not yet produce interpolated string nodes. This is planned as a future enhancement.

    \item \textbf{Closure method}: A generic \texttt{closure()} method for transitive navigation would replace hardcoded properties like \texttt{allSuperclasses} and \texttt{allSubclasses}.

    \item \textbf{\kw{forall} nesting}: The readability and semantics of nested \kw{forall} expressions need further validation.
\end{enumerate}


% ============================================
% 7. CONCLUSION
% ============================================
\section{Conclusion}\label{sec:conclusion}

This document has presented \jjtl{} and \jjel{} as a pair of complementary domain-specific languages for Model-Driven Engineering. The key contribution is \jjel{}'s set-theoretic approach to expressions: by reinterpreting \kw{forall} as a set comprehension rather than a logical assertion, we achieve a significant reduction in cognitive load while maintaining the declarative character valued in the MDE community.

The comparative analysis with ATL, EOL/ETL, QVT-R, and QVT-O reveals that each language occupies a distinct design point:

\begin{itemize}
    \item ATL combines declarative rules with OCL expressions but suffers from two-phase execution complexity and limited tooling.
    \item ETL prioritizes developer familiarity via an imperative style but sacrifices the declarative invertibility that enables advanced transformation analysis.
    \item QVT-R achieves full declarativity and bidirectionality at the cost of practical usability.
    \item QVT-O provides a pragmatic imperative approach with OCL's verbosity.
\end{itemize}

\jjtl{}/\jjel{} targets the space between ATL's declarative elegance and ETL's practical simplicity, aiming for a language that is \emph{declarative enough} for transformation analysis (invertibility, traceability) while being \emph{intuitive enough} for adoption by domain experts who may not have a formal logic background.


% ============================================
% APPENDIX
% ============================================
\appendix

\section{JjEL Grammar (Complete)}\label{app:grammar}

\begin{lstlisting}[style=grammar,caption={JjEL formal grammar in EBNF}]
expression     = forallExpr | existsExpr | withExpr | ifThenElse

forallExpr     = 'forall' IDENT 'in' expression
                   ('such' 'that' expression)?
                   (':' expression)?
                   ('do' expression)?

existsExpr     = 'exists' IDENT 'in' expression
                   (':' | 'such' 'that') expression

withExpr       = 'with' expression 'do' expression

ifThenElse     = 'if' expression 'then' expression ('else' expression)?
               | nullCoalesce

nullCoalesce   = implies ('??' implies)*

implies        = logicalOr ('implies' logicalOr)?

logicalOr      = logicalAnd ('or' logicalAnd)*

logicalAnd     = equality ('and' equality)*

equality       = comparison (('==' | '!=') comparison)*

comparison     = isType (('<' | '>' | '<=' | '>=') isType)*

isType         = addition ('is' IDENT)?

addition       = multiplication (('+' | '-') multiplication)*

multiplication = unary (('*' | '/' | '%') unary)*

unary          = 'not' unary
               | '-' unary
               | postfix

postfix        = primary (
                   '.' IDENT ('(' argList ')')?
                 | '?.' IDENT ('(' argList ')')?
                 | '[' expression ']'
                 )*

primary        = literal | IDENT | lambda | arrayLiteral
               | '(' expression ')'

lambda         = IDENT '=>' expression
               | '(' IDENT (',' IDENT)* ')' '=>' expression

arrayLiteral   = '[' (expression (',' expression)*)? ']'

argList        = (expression (',' expression)*)?

literal        = STRING | NUMBER | 'true' | 'false' | 'null'
\end{lstlisting}


\section{JjEL Built-in Methods (Complete)}\label{app:builtins}

\subsection{String Methods (33)}

\begin{longtable}{@{}lll@{}}
\toprule
\textbf{Method} & \textbf{Parameters} & \textbf{Description} \\
\midrule
\endhead
\texttt{toUpper()} & -- & Uppercase \\
\texttt{toLower()} & -- & Lowercase \\
\texttt{capitalize()} & -- & First letter uppercase \\
\texttt{uncapitalize()} & -- & First letter lowercase \\
\texttt{camelCase()} & -- & \texttt{"my name"} $\to$ \texttt{"myName"} \\
\texttt{pascalCase()} & -- & \texttt{"my name"} $\to$ \texttt{"MyName"} \\
\texttt{snakeCase()} & -- & \texttt{"myName"} $\to$ \texttt{"my\_name"} \\
\texttt{kebabCase()} & -- & \texttt{"myName"} $\to$ \texttt{"my-name"} \\
\texttt{trim()} & -- & Remove surrounding whitespace \\
\texttt{trimStart()} & -- & Remove leading whitespace \\
\texttt{trimEnd()} & -- & Remove trailing whitespace \\
\texttt{padStart(n, ch)} & length, char & Pad start to length \\
\texttt{padEnd(n, ch)} & length, char & Pad end to length \\
\texttt{contains(s)} & substring & Contains check \\
\texttt{startsWith(s)} & prefix & Prefix check \\
\texttt{endsWith(s)} & suffix & Suffix check \\
\texttt{indexOf(s)} & substring & First occurrence index \\
\texttt{lastIndexOf(s)} & substring & Last occurrence index \\
\texttt{matches(r)} & regex & Regex match \\
\texttt{replace(o, n)} & old, new & Replace first \\
\texttt{replaceAll(o, n)} & old, new & Replace all \\
\texttt{substring(s, e?)} & start, end? & Extract substring \\
\texttt{slice(s, e?)} & start, end? & Slice (negative indices) \\
\texttt{split(sep)} & separator & Split to array \\
\texttt{repeat(n)} & count & Repeat string \\
\texttt{reverse()} & -- & Reverse string \\
\texttt{length} & -- & Character count \\
\texttt{isEmpty} & -- & Length == 0 \\
\texttt{isNotEmpty} & -- & Length > 0 \\
\texttt{isBlank} & -- & Empty or whitespace only \\
\texttt{isNotBlank} & -- & Not blank \\
\texttt{charAt(i)} & index & Character at position \\
\texttt{toNumber()} & -- & Parse as number \\
\texttt{toInt()} & -- & Parse as integer \\
\bottomrule
\end{longtable}

\subsection{Collection Methods (28)}

After \kw{forall} replaces \texttt{filter()} and \texttt{map()}, the remaining collection methods are:

\begin{longtable}{@{}lll@{}}
\toprule
\textbf{Method} & \textbf{Parameters} & \textbf{Description} \\
\midrule
\endhead
\texttt{sum()} & -- & Numeric sum \\
\texttt{avg()} & -- & Numeric average \\
\texttt{min()} & -- & Minimum value \\
\texttt{max()} & -- & Maximum value \\
\texttt{count()} & -- & Element count \\
\texttt{first} & -- & First element \\
\texttt{last} & -- & Last element \\
\texttt{at(i)} & index & Element at index \\
\texttt{indexOf(x)} & item & Index of element \\
\texttt{size} & -- & Collection size \\
\texttt{length} & -- & Alias for size \\
\texttt{isEmpty} & -- & Size == 0 \\
\texttt{isNotEmpty} & -- & Size > 0 \\
\texttt{contains(x)} & item & Membership test \\
\texttt{distinct()} & -- & Remove duplicates \\
\texttt{distinctBy(fn)} & lambda & Remove duplicates by key \\
\texttt{sortBy(fn)} & lambda & Sort ascending \\
\texttt{sortByDesc(fn)} & lambda & Sort descending \\
\texttt{reverse()} & -- & Reverse order \\
\texttt{flatten()} & -- & Flatten nested arrays \\
\texttt{flatMap(fn)} & lambda & Map + flatten \\
\texttt{groupBy(fn)} & lambda & Group by key \\
\texttt{take(n)} & count & First $n$ elements \\
\texttt{skip(n)} & count & Skip first $n$ \\
\texttt{takeWhile(fn)} & predicate & Take while true \\
\texttt{skipWhile(fn)} & predicate & Skip while true \\
\texttt{join(sep)} & separator & Join to string \\
\texttt{all(fn)} & predicate & All satisfy? \\
\texttt{any(fn)} & predicate & Any satisfies? \\
\texttt{none(fn)} & predicate & None satisfy? \\
\bottomrule
\end{longtable}


\section{JjTL Token Types}\label{app:tokens}

\begin{center}
\begin{tabular}{@{}lll@{}}
\toprule
\textbf{Category} & \textbf{Token} & \textbf{Example} \\
\midrule
\multirow{5}{*}{Keywords} & \texttt{TRANSFORMATION} & \texttt{transformation} \\
& \texttt{FROM} & \texttt{from} \\
& \texttt{TO} & \texttt{to} \\
& \texttt{WHEN} & \texttt{when} \\
& \texttt{HELPER} & \texttt{helper} \\
\midrule
\multirow{4}{*}{Interactive} & \texttt{ALERT} & \texttt{alert} \\
& \texttt{NOTIFY} & \texttt{notify} \\
& \texttt{PROMPT} & \texttt{prompt} \\
& \texttt{INPUT} & \texttt{input} \\
\midrule
\multirow{4}{*}{Literals} & \texttt{STRING} & \texttt{"hello"} \\
& \texttt{NUMBER} & \texttt{42}, \texttt{3.14} \\
& \texttt{BOOLEAN} & \texttt{true}, \texttt{false} \\
& \texttt{IDENTIFIER} & \texttt{name}, \texttt{State} \\
\midrule
\multirow{3}{*}{Operators} & \texttt{ARROW} & \texttt{->} \\
& \texttt{DOT} & \texttt{.} \\
& \texttt{QUESTION\_DOT} & \texttt{?.} \\
\bottomrule
\end{tabular}
\end{center}


\end{document}
